% Pass 2: teaching example derived from Veach-Guibas and Veach MIS lineage.
\clearpage
\section*{Research reconstruction: combine two samplers without double-counting}
\noindent\textbf{Primary sources:}
Eric Veach and Leonidas Guibas,
\href{https://lightfield.stanford.edu/papers/combine/}
{\emph{Optimally Combining Sampling Techniques for Monte Carlo Rendering}
(SIGGRAPH 1995)}, and Veach's
\href{https://graphics.stanford.edu/papers/veach_thesis/}
{1997 dissertation, Chapter 9 on multiple importance sampling}.
This is a one-dimensional \emph{balance-heuristic}
laboratory, not a replication of the authors' light-transport
images or their efficiency figures.

\paragraph{A quantity with an exact answer.}
Take
\[
I=\int_0^1 e^x\,dx=e-1.
\]
A uniform sampler uses \(p_1(x)=1\) on \([0,1]\).
A second sampler uses \(p_2(x)=x+1/2\), which is also
a normalized, strictly positive density:
\(\int_0^1(x+1/2)\,dx=1\).
The second proposal samples larger \(x\) values more
frequently, where \(e^x\) is also larger. Neither proposal
changes the mathematical integral being estimated.

A sample from \(p_2\) can be generated by inverting its
CDF \(F_2(x)=(x+x^2)/2\). For a uniform
\(u\in[0,1)\), set
\[
x=\frac{\sqrt{1+8u}-1}{2}.
\]
This transformation is worth deriving: if the implemented
inverse CDF does not match the claimed PDF, subsequent
importance weights are invalid even when the images look smooth.

\paragraph{The multi-sample balance heuristic.}
Draw \(n\) samples from each proposal.
With equal allocation, define the mixture density
\[
m(x)=\frac{p_1(x)+p_2(x)}2=\frac34+\frac x2.
\]
An unbiased estimator is
\[
\widehat I_{\rm balance}
=\frac1{2n}\left[
\sum_{i=1}^n\frac{e^{X_i}}{m(X_i)}
+\sum_{j=1}^n\frac{e^{Y_j}}{m(Y_j)}
\right],\qquad
X_i\sim p_1,\quad Y_j\sim p_2.
\]
Take its expectation under independent draws:
\[
\mathbb E[\widehat I_{\rm balance}]
=\frac12\int_0^1\frac{e^x p_1(x)}{m(x)}\,dx
+\frac12\int_0^1\frac{e^x p_2(x)}{m(x)}\,dx
=\int_0^1 e^x\,dx=I.
\]
The last equality follows because
\(\frac12[p_1(x)+p_2(x)]=m(x)\).
This establishes unbiasedness under the declared normalized
proposals and valid support, not low variance for every
integrand or every proposal choice.

\paragraph{A finite-sample experiment with honest inference.}
The independent companion
\texttt{tools/labs/book02\_papers\_sampling.py}
runs three estimators: uniform importance sampling,
tilted-proposal importance sampling, and the balance
heuristic above. It records the mean and sample variance
over \emph{independent seeded replications}.
If one technique happens to dominate on a particular
simple integrand, the balance heuristic may have
higher measured variance than that specialized technique.
This is not a contradiction: the motivation for MIS is
robustly combining useful strategies when no single
proposal is best over the entire domain.

\paragraph{Failure conditions.}
The identity requires the mixture density to be
positive wherever the integrand contributes.
A proposal that never samples a region of nonzero target
mass cannot recover that region by assigning arbitrarily
large weights to the sampled data. Mislabeling an area
PDF as a solid-angle PDF also invalidates the estimator;
different measures require a geometric Jacobian.
Reusing the same seed across all experimental
replications makes identical pseudo-observations,
not independent inferential evidence.

\paragraph{Research extension.}
Compare a balanced multi-sample estimator, a
single-mixture estimator, stratified sampling and
the power heuristic. Measure variance against
\emph{actual total sample cost}, not only per-technique
sample count, and disclose whether the samples are
independent and which expectations remain unbiased.
Full path-tracing integration belongs in Book Four;
this volume develops the mathematical contracts.
